\chapter*{Abstract}

\thispagestyle{front}

{
    \color{gray}
    Federated learning is a machine learning approach that allows multiple participants to collaboratively train a model without sharing their local data. This approach is particularly useful in scenarios where data privacy and security are a major concern. But the cybersecurity landscape conducted researches and studies, highlighting data poisoning attacks as one of the main concern in federated learning environments. 
    Securing the models performance is a concern in which the aggregator node became responsible for by implementing a robust methodogoly against data poisoning attacks. 
    
    In this paper, a web-based open-source simulation platform, called DPArena, has been developed to address the need of finding the most robust federated learning configuration preventing model degradation and generalization loss. 
    Each simulation has its own configuration providing data scientist with a dedicated environment for testing its machine learning algorithm againts different data poisoning attacks. The paper ends with two case studies demonstrating the effectiveness of the platform.
}

\thispagestyle{front}

\newpage

\chapter*{Rezumat}

\thispagestyle{front}

{
    \color{gray}
    
    Învățarea federată este o abordare a învățării automate care permite mai multor participanți să antreneze în mod colaborativ un model fără a-și partaja datele locale. Această abordare este utilă în special în scenariile în care confidențialitatea și securitatea datelor reprezintă o preocupare majoră. Însă, în domeniul securității cibernetice s-au efectuat cercetări și studii, evidențiind atacurile de otrăvire a datelor ca fiind una dintre principalele preocupări în mediile de învățare federate.
    Securizarea performanței modelelor este o preocupare pentru care nodul agregator a devenit responsabil prin implementarea unei metodologii robuste împotriva atacurilor de otrăvire a datelor.


    În această lucrare, a fost dezvoltată o platformă web de simulare open-source, numită DPArena, pentru a răspunde nevoii de a găsi cea mai robustă configurație de învățare federată, prevenind degradarea modelului și pierderea generalizării.
    Fiecare simulare are propria configurație, oferind specialiștilor în date un mediu dedicat pentru testarea algoritmului său de învățare automată împotriva diferitelor atacuri de otrăvire a datelor.
    Lucrarea se încheie cu două studii de caz care demonstrează eficacitatea platformei.
}

\thispagestyle{front}
